\chapter*{Streszczenie}
Niniejsza praca poświęcona jest analizie możliwości współczesnych modeli automatycznego rozpoznawania mowy (ASR) w zakresie uczenia semantycznych reprezentacji danych, zarówno w ujęciu wielojęzycznym, jak i w odniesieniu do słownictwa specjalistycznego. W pierwszym etapie przeanalizowano reprezentacje przestrzeni ukrytej różnych modeli ASR, oceniając stopień, w jakim odzwierciedlają one informacje semantyczne. W tym celu wykorzystano reprezentacje wektorowe transkrypcji uzyskane za pomocą trzech różnych modeli reprezentacji wektorowej (tj. embeddingowych) tekstu. Następnie zaproponowano dwie nowe architektury głębokich sieci neuronowych, wykorzystujące istniejące modele ASR i przeznaczone do multimodalnego wyszukiwania dokumentów na podstawie zapytań w postaci sygnału mowy. W ostatnim etapie wyniki wcześniejszych analiz wykorzystano do opracowania kompleksowego systemu wyszukiwania dokumentów dla specjalistycznych danych medycznych w języku polskim. System ten posłużył następnie do porównania i oceny różnych rozwiązań stosowanych w realizacji tego zadania.
\chapter*{Abstract}
This work investigates the capabilities of contemporary automatic speech recognition (ASR) models to learn semantic representations of data, both in multilingual settings and in the context of specialized vocabulary. First, latent-space representations derived from different ASR models were analyzed to assess the extent to which they capture semantic information. For this purpose, vector representations of the corresponding transcriptions, generated using three different text embedding models, were employed. Next, two novel deep-learning architectures based on existing ASR models were proposed for multimodal document retrieval using spoken queries. Finally, the findings from the preceding analyses were used to develop a comprehensive document retrieval system for specialized Polish medical data. The resulting system was subsequently employed to compare and evaluate different approaches to this task.
